\section{A canonical detected quotient of class four}
\label{sec:canonical-class-four}

Retain the canonical quotients $Q_2$ and $Q_3$, and their
source kernels $K_2$ and $K_3$. Put $B=\ker(H\to W)$.
Define
\[
 J_4=\overline{[H,K_2][B,B]},\qquad
 J_5=\overline{[H,K_3][B,K_2]},\qquad
 R_4=K_3/J_5,\qquad S_4=J_4/J_5.
\]
Set
\[
 \mathcal R=K_2/J_4,\qquad R_3=K_2/K_3,\qquad
 \bar E=B/K_2,\qquad N=\widehat\bigwedge^2W.
\]
All commutator subgroups and tensor products in this
section are closed and compact, respectively.



\begin{theorem}\label{thm:canonical-class-four}
For every odd prime $p$, the $\Lambda$-torsion submodule
$(R_4)_{\rm tor}$ is closed. Its inverse image
$K_4^{\rm det}$ in $K_3$ is fully Galois-stable and
contains the kernel of the actual outer action. The
canonical quotient
\[
                         Q_4=H/K_4^{\rm det}
\]
is torsionfree of nilpotency class exactly four and
fits into a central extension
\[
 1\longrightarrow R_4/(R_4)_{\rm tor}
 \longrightarrow Q_4\longrightarrow Q_3\longrightarrow1.
\]
The natural map $\operatorname{Lie}^4_{\mathbf Z_p}(W)
\to R_4/(R_4)_{\rm tor}$ is injective. In particular,
\[
 \Omega_p^{K_3}\subsetneq\Omega_p^{K_4^{\rm det}}
                              \subseteq\Omega_p^*.
\]
The quotient is the compact image of $H$ in the full
rational outer image through generator length four.
\end{theorem}

\begin{proof}[Proof of Theorem~\ref{thm:canonical-class-four}]
We first calculate the source module integrally, then
prove its rational coefficient injection, and finally
identify its torsion by one fixed annihilator.

The established filtration
$\mathcal F_1=H,\mathcal F_2=B,\mathcal F_3=K_2,\mathcal F_4=K_3$ satisfies
$[\mathcal F_1,\mathcal F_2]\subset \mathcal F_3$ and
$[\mathcal F_1,\mathcal F_3][\mathcal F_2,\mathcal F_2]\subset \mathcal F_4$.
The three-subgroup identity gives
$[H,[B,B]]\subseteq\overline{[B,K_2]}$.
It follows that $J_5\subseteq J_4$, and that $R_4,S_4$
are central compact $\mathbf Z_p$-modules in $H/J_5$.
Their arithmetic actions factor through
$\Gamma\times\Delta$. The exact quotient is
\begin{equation}\label{eq:class4-old-torsion}
 0\longrightarrow S_4\longrightarrow R_4
     \longrightarrow \mathcal R_{\rm tor}\longrightarrow0.
\end{equation}
Here $\mathcal R_{\rm tor}$ is the entire old torsion module,
including its possibly infinite bounded $p$-torsion.

Let $\nu:\widehat\bigwedge^2W\to\bar E$ and
$\beta:W\widehat\otimes\bar E\to R_3$ be the established
bracket maps. Define
\[
 A_4=(W\widehat\otimes R_3)\oplus
                          \widehat\bigwedge^2\bar E
\]
and let $I_{112}$ be the image of
\begin{align*}
 j_{112}:\widehat\bigwedge^2W\widehat\otimes\bar E
          &\longrightarrow A_4,\\
 (x\wedge y)\otimes e&\longmapsto
 \bigl(x\otimes\beta(y,e)-y\otimes\beta(x,e),
                         -\nu(x\wedge y)\wedge e\bigr).
\end{align*}
The exact integral presentation is
\begin{equation}\label{eq:class4-presentation}
                         S_4=A_4/I_{112}.
\end{equation}
Indeed, the first summand maps by the commutator of an
$H$-lift with a $K_2$-lift. Changing the lifts contributes
$[B,K_2]$ or $[H,K_3]$. The second summand maps by the
commutator of two $B$-lifts; changing a lift contributes
$[B,K_2]$. Centrality gives bilinearity and alternation,
and Hall--Witt gives the displayed relation.

We prove exactness by integral continuous homology.
Write $B_3=B/K_3$. The established class-three structure
gives an abelian group
$B_3=\bar E^+\oplus R_3^-$, with $R_3$ central, and
the $W$-action is
\[
 w\cdot e=e+\beta(w,e),\qquad w\cdot r=r.
\]
The Hopf formula for the free pro-$p$ presentation
$H\twoheadrightarrow Q_3$ identifies
\[
 H_2(Q_3,\mathbf Z_p)
             =(K_3\cap[H,H])/[H,K_3].
\]
Its map to $K_3/J_4=\mathcal R_{\rm tor}$ has image
$\ker(\mathcal R_{\rm tor}\to H^{\rm ab})$. This image is odd:
the kernel lies in the old commutator module $\mathcal S$,
which is odd. Therefore the even part of $H_2(Q_3)$
is represented entirely by $J_4$. The commutator
classes of $[B,K_2]$ are the image of the commuting-pair
map $B_3\widehat\wedge R_3\to H_2(Q_3)$.
Consequently the actual source quotient is
\begin{equation}\label{eq:class4-Hopf}
 S_4=\bigl(H_2(Q_3,\mathbf Z_p)/
             \operatorname{im}(B_3\widehat\wedge R_3)\bigr)^+.
\end{equation}

Apply Hochschild--Serre homology to
$1\to B_3\to Q_3\to W\to1$.
Put $U_4=W\widehat\otimes R_3$,
$V_4=\widehat\bigwedge^2\bar E$ and
$P_4=\widehat\bigwedge^2W\widehat\otimes\bar E$.
Write $\delta:P_4\to U_4$ and $\epsilon:P_4\to V_4$
for the two components of $j_{112}=(\delta,-\epsilon)$.
Use the completed Koszul complex with symmetric
parameters $z_i=w_i^{1/2}-w_i^{-1/2}$ for a compact
$\mathbf Z_p$-basis of $W$. They differ from $w_i-1$
by units and satisfy $c(z_i)=-z_i$. On $B_3$ the
parameter acts as $\beta_i$: if $w_i=1+\beta_i$ on
this module, then $\beta_i^2=0$ and the two half powers
give $z_i=\beta_i$. Thus its chain parity is literally
the exterior-degree parity times the coefficient
parity. The complex gives
\[
 E^2_{1,1}{}^+=U_4/\operatorname{im}\delta,
 \qquad E^2_{2,1}{}^+=\ker\delta.
\]
Indeed, every product of two augmentation operators
on $B_3$ is zero. The even one-chains are the
$W\widehat\otimes R_3$ cycles, and their boundaries
are exactly $\delta(P_4)$. The even two-chains identify
with $P_4$, with differential $\delta$; even three-chain
boundaries vanish because their coefficients lie in
the trivial $W$-module $R_3$.

After quotienting by the image of
$B_3\widehat\wedge R_3$, the even $(0,2)$ term is
$V_4$. Every $W$-coinvariant relation in
$\widehat\bigwedge^2B_3$ contains an $R_3$ factor,
so no further relation is introduced in $V_4$.
The transgression from $\ker\delta$ is $-\epsilon$;
on lifted commuting pairs this is exactly the
Hall--Witt formula already displayed. The only later
possible incoming differential is $d_3$ from a
subquotient of $\widehat\bigwedge^3W$, which is odd;
its even target makes it zero. Likewise the only
possible incoming differential to $E^2_{1,1}{}^+$
has odd source and is zero. Finally the transgression
on $E^2_{2,0}=\widehat\bigwedge^2W$ is $\nu$, whose
even target is $\bar E$ and which is injective.
Thus no $(2,0)$ piece survives.

It follows that \eqref{eq:class4-Hopf} has the exact
two-step filtration
\[
 0\longrightarrow V_4/\epsilon(\ker\delta)
 \longrightarrow S_4
 \longrightarrow U_4/\operatorname{im}\delta
 \longrightarrow0.
\]
The canonical surjection $A_4/I_{112}\to S_4$ has
exactly this filtration by the image of $V_4$.
Its maps on both pieces are the identity, represented
respectively by two $\bar E$ lifts and by a $W$ lift
with an $R_3$ lift. It is therefore an isomorphism.
The completed Koszul calculations follow from finite
coordinate complexes and their compact closed images;
the image of $j_{112}$ is itself compact and closed.
This proves \eqref{eq:class4-presentation} without
assuming primitivity of $N\subset\bar E$ and without
dividing by $3$, so it includes $p=3$.

We now give a uniform integral bound. Choose $b$ such
that $p^b(T/\pi Z)[p^\infty]=0$, as in the class-three
proof, and put
\[
 D_0=\operatorname{im}(\mathcal R_{\rm tor}^-\to D).
\]
Choose $c$ such that $p^c(D/D_0)[p^\infty]=0$.
The established exact sequence is
\[
 0\longrightarrow S_3^{\rm tf}\longrightarrow R_3
 \longrightarrow D/D_0\longrightarrow0,
 \qquad S_3^{\rm tf}=\mathcal S/\mathcal S[p^\infty].
\]
Choose a finite free $\mathbf Z_p$-module $C_2\subset\bar E$
lifting a direct summand complementary to the torsion
of $T/\pi Z$. Then
\[
 E_0=N\oplus C_2\subset\bar E,\qquad p^b\bar E\subset E_0.
\]
The free weighted Hall module
\[
 L_3(W,C_2)=\operatorname{Lie}^3(W)
                      \oplus(W\widehat\otimes C_2)
\]
injects into $S_3^{\rm tf}$, and its cokernel is killed
by $p^b$. To see injectivity, the Jacobi image is already
in $W\widehat\otimes N$; the resulting free Hall module
is torsionfree and meets the removed $p$-torsion trivially.
Choose a finite free $C_3\subset R_3$ lifting a complement
to the torsion of $D/D_0$. Thus
\[
 R_0=L_3(W,C_2)\oplus C_3\subset R_3,
 \qquad p^{b+c}R_3\subset R_0.
\]
These auxiliary lattices need not be $\Lambda$-stable.

Put
\[
 A_0=(W\widehat\otimes R_0)\oplus
                          \widehat\bigwedge^2E_0,
 \qquad d=\max\{b+c,2b\}.
\]
Then $A_0\subset A_4$ and $p^dA_4\subset A_0$.
The tensor inclusions are injective: compact torsionfree
$\mathbf Z_p$-modules are products of copies of
$\mathbf Z_p$, and tensoring in these coordinates
preserves the displayed inclusions. Let $I_0$ be the
Jacobi image with the third input in $E_0$. It lies in
$A_0$, and
\[
 A_0/I_0=L_4(W,C_2,C_3),\qquad p^bI_{112}\subset I_0,
\]
where $L_4$ is degree four of the free graded Lie algebra
on generators of weights $1,2,3$, respectively.
For the first equality, its lower pieces are the free
weighted pieces already displayed. The only new Jacobi
identities in degree four have type $(1,1,2)$, exactly
$I_0$. This is the universal free graded Lie algebra
truncated at degree four; a Hall basis proves its
integral torsionfreeness.

If $a\in A_4$ represents a $p$-torsion element of $S_4$,
then $p^ka\in I_{112}$ for some $k$. Since
$p^da\in A_0$ and $p^{k+b}p^da\in I_0$, saturation
of $I_0$ in $A_0$ implies $p^da\in I_0$. Therefore
\begin{equation}\label{eq:class4-pbound}
                           p^dS_4[p^\infty]=0.
\end{equation}
Furthermore $A_0/I_0\to S_4$ is injective: its kernel
is killed by $p^b$, while $A_0/I_0$ is torsionfree.
In particular
\begin{equation}\label{eq:class4-pure-Hall}
                  \operatorname{Lie}^4(W)\hookrightarrow S_4.
\end{equation}

It remains to prove rational detection of $S_4$.
Work over a finite splitting field $L/\mathbf Q_p$ and
write
\[
 T'=(T/\pi Z)_L,\qquad D'=(D/D_0)_L.
\]
Auxiliary $L$-linear sections give
\[
 \bar E_L=N_L\oplus T',\qquad
 R_{3,L}=L_3(W_L,T')\oplus D'.
\]
Consequently \eqref{eq:class4-presentation} gives
\begin{equation}\label{eq:class4-weighted-free}
                        S_{4,L}=L_4(W_L,T',D'),
\end{equation}
with weights $1,2,3$. These sections are used only for
the proof; the source quotient remains intrinsic.

The target throughout is the free Lie algebra of the
full rational outer image, with its own generator
length. Choose the established finite set of odd
$\Lambda[1/p]$-colors for $W[1/p]$. Coefficient arrays
are bounded power series in independent position
variables $z_i=1+X_i$, with the companion denominators
below inverted.

For a Jordan block of $T'$ with eigenvalue $\alpha$
and length $r$, take a top lift $y\in\bar E_L$ and put
$C=(\gamma-\alpha)^ry\in N_L$. The companions of
distinct blocks at the same root have independent
constant-field residues modulo $\gamma-\alpha$.
Otherwise a linear relation would produce a nonzero
$\Lambda_L$-torsion element in $\bar E_L$, whose image
is a nonzero combination of the independent Jordan
socles. The ordered pair kernels of the successive
Jordan vectors are therefore
\[
              \frac{C(z_1,z_2)}{(z_1z_2-\alpha)^{r-j}},
              \qquad 0\le j<r.
\]
Similarly, for a Jordan block of $D'$ with eigenvalue
$\beta$ and length $s$, a top lift $v\in R_{3,L}$ has
companion $C'=(\gamma-\beta)^sv\in S_{3,L}$.
The $\Lambda_L$-torsionfreeness of $R_{3,L}$ gives
independent constant-field companion residues.
Proposition~\ref{prop:degree-three-saturation} transports
this independence to their actual coefficient residues
on $z_1z_2z_3=\beta$. The ordered triple kernels are
\[
        \frac{C'(z_1,z_2,z_3)}{(z_1z_2z_3-\beta)^{s-j}},
                         \qquad0\le j<s.
\]
The numerator may have internal pair poles, which are
coprime to the displayed triple divisor.
Lemma~\ref{lem:positive-eigenvalues} ensures that these
denominators do not vanish anywhere on the positive
odd-weight moment grid. Thus the formulas apply to
the actual target without discarding any odd generator.

Here is an explicit injectivity argument, which also
avoids any assumption about total degree-four
divisibility. Embed the free weighted Lie algebra
in its free associative tensor algebra and induct on
total degree. Every source word starts with a $W$,
$T'$ or $D'$ token. At the prefix divisor
$z_1z_2z_3-\beta$, only an initial $D'$ token can have
a pole. A later triple token has a different divisor;
all internal pair divisors are also distinct. At the
largest pole order the residue is a finite sum
\[
                   \sum_i\overline C'_i\,G_i,
\]
where the $G_i$ depend only on the disjoint tail
variables. The finite list $\overline C'_i$ is
independent over $L$. Clear its internal pair
denominators and take finitely many color and
coefficient functionals giving an invertible matrix
over $L$. Applying this matrix coefficientwise to the
tail variables forces every $G_i$ to vanish. The
induction hypothesis makes each corresponding source
tail zero. Descend through the finite list of pole
orders and roots to remove all initial $D'$ terms.

Next use the prefix divisor $z_1z_2-\alpha$. Only an
initial $T'$ token can now have this pole, and the same
residue argument removes all such terms. The remaining
terms begin with $W$ and are detected by the ordinary
first-slot coefficient map tensored with the already
injective tail map. The induction starts in degree
one. It proves injectivity in every fixed degree of
this auxiliary free weighted algebra, in particular
in degree four of \eqref{eq:class4-weighted-free}.
The argument extends residue independence only to
disjoint tail variables, never to the function field
of the variables inside the companion.

A zero actual target coefficient first gives zero on
the odd moment grid. Clearing the finitely many
companion denominators and using bounded power-series
uniqueness in each variable gives the coefficient
identity just used. We have proved
\begin{equation}\label{eq:class4-rational-injection}
 S_4[1/p]\hookrightarrow
                     \text{target generator length four}.
\end{equation}
It also follows that $S_4[1/p]$ is $\Lambda$-torsionfree:
a nonzero $f(\gamma)\in\Lambda$ acts on each ordered
coefficient array by the nonzero multiplier
$f(z_1z_2z_3z_4)$ in a domain.

The preceding class-three construction gives
$\mathcal R_{\rm tor}=\ker(p^ba(\gamma):\mathcal R\to\mathcal R)$, where $a$
is the characteristic polynomial on $D\oplus\pi Z$.
By \eqref{eq:class4-old-torsion},
$p^ba(\gamma)R_4\subseteq S_4$. Equations
\eqref{eq:class4-pbound} and
\eqref{eq:class4-rational-injection} now give the exact
closed-submodule formula
\begin{equation}\label{eq:class4-closed-torsion}
 (R_4)_{\rm tor}
       =\ker\bigl(p^{b+d}a(\gamma):R_4\to R_4\bigr).
\end{equation}
Indeed, the $p^ba(\gamma)$-multiple of a torsion element
lies in $S_4$ and is rationally zero, hence is killed
by $p^d$. The converse is immediate from the displayed
nonzero annihilator. The formula proves closedness;
the torsion property and the source construction prove
full arithmetic invariance.

The exact class-three target kernel is $K_3$:
the prior detection argument identifies its possible
central ambiguity as old $\Lambda$-torsion, and
Lemma~\ref{lem:positive-eigenvalues} makes that torsion
zero on every actual positive-weight coordinate.
Consequently $J_5$ has target length at least five,
and the map on $R_4$ takes values in pure length four.
An actual outer-kernel element has image $r\in R_4$
with zero coefficient image. Its multiple
$p^ba(\gamma)r\in S_4$ is zero rationally by
\eqref{eq:class4-rational-injection}, and
\eqref{eq:class4-pbound} gives
$p^{b+d}a(\gamma)r=0$. Formula
\eqref{eq:class4-closed-torsion} proves the asserted
kernel containment.

Conversely the roots of $a$ are eigenvalues of
$D\oplus T$, so Lemma~\ref{lem:positive-eigenvalues}
shows that $a(\gamma)$ is nonzero on every positive
target coordinate of length four. Thus
$(R_4)_{\rm tor}$ itself has zero target image.
This proves the exact compact-image assertion.
The displayed central extension is torsionfree since
$Q_3$ is torsionfree and the central kernel has no
$p$-torsion. The pure Hall submodule
\eqref{eq:class4-pure-Hall} is $\Lambda$-torsionfree by
its four-variable coefficient injection, so it meets
$(R_4)_{\rm tor}$ trivially. Finally
$\gamma_4H\subseteq K_3$ and $\gamma_5H\subseteq J_5$,
while $[x,[x,[x,y]]]$ is a nonzero surviving Hall word
for independent odd coordinates $x,y$. This gives
class exactly four and strictness over $Q_3$.
\end{proof}

The exact abelianization records any even directions
restored at this stage. Since $S_4$ is even, the odd
part of $R_4$ is entirely $\Lambda$-torsion and maps
onto $D_0$ in $H^{\rm ab}$. Put
\[
 Z_0=\operatorname{im}\bigl((R_4)_{\rm tor}^+
                                     \longrightarrow\pi Z\bigr).
\]
Then the abelian image of $K_4^{\rm det}$ is exactly
$D_0\oplus Z_0$, and hence
\[
                       Q_4^{\rm ab}=(U/D_0)\oplus(T/Z_0).
\]
This identity is induced by the actual abelianization
map of the source; no vanishing of $Z_0$ is asserted.
